%!TEX program = xelatex
\documentclass[12pt,a4paper,oneside]{article}

% ============================================================================
% PACKAGES & CONFIGURATION
% ============================================================================
\usepackage[utf-8]{inputenc}
\usepackage[left=2.5cm,right=2.5cm,top=3cm,bottom=3cm]{geometry}
\usepackage{graphicx}
\usepackage{fancyhdr}
\usepackage{color}
\usepackage{listings}
\usepackage{hyperref}
\usepackage{xcolor}
\usepackage{tcolorbox}
\usepackage{array}
\usepackage{booktabs}
\usepackage{amsmath}
\usepackage{amssymb}
\usepackage{float}
\usepackage{titlesec}
\usepackage{enumitem}
\usepackage{fancyvrb}

% Color definitions
\definecolor{codegreen}{rgb}{0.0, 0.5, 0.0}
\definecolor{codegray}{rgb}{0.5, 0.5, 0.5}
\definecolor{codepurple}{rgb}{0.58, 0, 0.82}
\definecolor{backcolour}{rgb}{0.95, 0.95, 0.95}
\definecolor{lightblue}{rgb}{0.9, 0.95, 1.0}
\definecolor{darkblue}{rgb}{0.0, 0.2, 0.6}

% Code listing configuration
\lstset{
    backgroundcolor=\color{backcolour},
    commentstyle=\color{codegreen},
    keywordstyle=\color{codepurple},
    numberstyle=\tiny\color{codegray},
    stringstyle=\color{codegreen},
    basicstyle=\ttfamily\small,
    breakatwhitespace=false,
    breaklines=true,
    captionpos=b,
    keepspaces=true,
    numbers=left,
    numbersep=5pt,
    showspaces=false,
    showstringspaces=false,
    showtabs=false,
    tabsize=2,
    frame=single,
    language=json,
}

% Header and footer
\pagestyle{fancy}
\fancyhf{}
\rhead{\textcolor{darkblue}{\textit{VisioLearn API Documentation}}}
\lhead{\textcolor{darkblue}{v1.0}}
\cfoot{\thepage}
\renewcommand{\headrulewidth}{0.5pt}
\renewcommand{\footrulewidth}{0.5pt}

% Section styling
\titleformat{\section}{\Large\bfseries\color{darkblue}}{}{0pt}{}[\titlerule]
\titleformat{\subsection}{\large\bfseries\color{darkblue}}{}{0pt}{}
\titleformat{\subsubsection}{\bfseries\color{darkblue}}{}{0pt}{}

% ============================================================================
% TITLE PAGE
% ============================================================================
\title{%
    \vspace{2cm}
    \Huge\textbf{\textcolor{darkblue}{VisioLearn Backend API}} \\
    \vspace{0.5cm}
    \Large\textit{Comprehensive API Documentation} \\
    \vspace{2cm}
    \large\textbf{AI-Powered Audio Learning Platform}
}

\author{%
    VisioLearn Development Team \\
    \vspace{0.5cm}
    \href{https://visiolearn-backend.onrender.com}{https://visiolearn-backend.onrender.com}
}

\date{\vspace{1cm}\large \textbf{Version 1.0} \\ May 2026}

% ============================================================================
% DOCUMENT
% ============================================================================
\begin{document}

% Title page
\maketitle
\thispagestyle{empty}
\newpage

% ============================================================================
% TABLE OF CONTENTS
% ============================================================================
\tableofcontents
\newpage

% ============================================================================
% SECTION 1: INTRODUCTION
% ============================================================================
\section{Introduction}

VisioLearn is an AI-powered audio learning platform designed for educational institutions. The backend API enables:

\begin{itemize}
    \item \textbf{Content Management}: Upload and process lesson content (PDF, DOCX, TXT)
    \item \textbf{Authentication}: Secure JWT-based user authentication with token rotation
    \item \textbf{Voice Sessions}: Track student learning interactions through voice commands
    \item \textbf{User Management}: Manage teachers, students, and administrative users
    \item \textbf{Analytics}: Track student progress and learning analytics
\end{itemize}

\subsection{API Base URL}
\begin{tcolorbox}[colback=lightblue, colframe=darkblue]
\texttt{https://visiolearn-backend.onrender.com/api/v1}
\end{tcolorbox}

\subsection{API Documentation}
Interactive API documentation is available at:
\begin{itemize}
    \item \textbf{Swagger UI}: \url{https://visiolearn-backend.onrender.com/docs}
    \item \textbf{ReDoc}: \url{https://visiolearn-backend.onrender.com/redoc}
    \item \textbf{OpenAPI Schema}: \url{https://visiolearn-backend.onrender.com/openapi.json}
\end{itemize}

% ============================================================================
% SECTION 2: AUTHENTICATION
% ============================================================================
\section{Authentication}

All API endpoints (except login) require authentication via JWT bearer token.

\subsection{Token Types}

VisioLearn uses two types of tokens:

\begin{itemize}
    \item \textbf{Access Token}: Short-lived token (60 minutes) for API requests
    \item \textbf{Refresh Token}: Long-lived token (90 days) for obtaining new access tokens
\end{itemize}

\subsection{Authorization Header Format}

All authenticated requests must include:

\begin{lstlisting}[language=bash]
Authorization: Bearer <access_token>
\end{lstlisting}

\subsection{User Roles}

VisioLearn supports three user roles:

\begin{table}[H]
\centering
\begin{tabular}{|l|l|}
\hline
\textbf{Role} & \textbf{Description} \\
\hline
\texttt{admin} & Full access to all endpoints and resources \\
\hline
\texttt{teacher} & Can upload lessons, view student progress \\
\hline
\texttt{student} & Can access lessons, participate in voice sessions \\
\hline
\end{tabular}
\caption{User Roles in VisioLearn}
\end{table}

% ============================================================================
% SECTION 3: API ENDPOINTS
% ============================================================================
\section{API Endpoints}

\subsection{Authentication Endpoints}

\subsubsection{1. Login}

\textbf{POST} \texttt{/auth/login}

Authenticate a user and receive access and refresh tokens.

\textbf{Request Body:}
\begin{lstlisting}
{
  "email": "admin@visiolearn.org",
  "password": "SecurePass123!"
}
\end{lstlisting}

\textbf{Response (200 OK):}
\begin{lstlisting}
{
  "access_token": "eyJhbGciOiJIUzI1NiIsInR5cCI6IkpXVCJ9...",
  "refresh_token": "eyJhbGciOiJIUzI1NiIsInR5cCI6IkpXVCJ9...",
  "token_type": "bearer"
}
\end{lstlisting}

\textbf{Error Responses:}
\begin{table}[H]
\centering
\begin{tabular}{|l|l|}
\hline
\textbf{Status} & \textbf{Reason} \\
\hline
401 & Incorrect email or password \\
403 & Account is deactivated \\
\hline
\end{tabular}
\end{table}

\subsubsection{2. Refresh Token}

\textbf{POST} \texttt{/auth/refresh}

Obtain a new access token using a valid refresh token.

\textbf{Request Body:}
\begin{lstlisting}
{
  "refresh_token": "eyJhbGciOiJIUzI1NiIsInR5cCI6IkpXVCJ9..."
}
\end{lstlisting}

\textbf{Response (200 OK):}
\begin{lstlisting}
{
  "access_token": "eyJhbGciOiJIUzI1NiIsInR5cCI6IkpXVCJ9...",
  "refresh_token": "eyJhbGciOiJIUzI1NiIsInR5cCI6IkpXVCJ9...",
  "token_type": "bearer"
}
\end{lstlisting}

\textbf{Notes:}
\begin{itemize}
    \item Old refresh token is automatically revoked (token rotation)
    \item Returns new access token and refresh token
    \item If refresh token is expired or revoked, returns 401
\end{itemize}

\subsubsection{3. Logout}

\textbf{POST} \texttt{/auth/logout}

Revoke a refresh token and log out the user.

\textbf{Headers:}
\begin{lstlisting}
Authorization: Bearer <access_token>
\end{lstlisting}

\textbf{Request Body:}
\begin{lstlisting}
{
  "refresh_token": "eyJhbGciOiJIUzI1NiIsInR5cCI6IkpXVCJ9..."
}
\end{lstlisting}

\textbf{Response (200 OK):}
\begin{lstlisting}
{
  "success": true,
  "message": "Successfully logged out"
}
\end{lstlisting}

% ============================================================================
% SECTION 4: LESSON NOTES ENDPOINTS
% ============================================================================
\subsection{Lesson Notes Endpoints}

\subsubsection{1. Upload Lesson Note}

\textbf{POST} \texttt{/notes/upload}

Upload a lesson document for processing and content extraction.

\textbf{Headers:}
\begin{lstlisting}
Authorization: Bearer <access_token>
Content-Type: multipart/form-data
\end{lstlisting}

\textbf{Form Parameters:}
\begin{table}[H]
\centering
\begin{tabular}{|l|l|l|}
\hline
\textbf{Parameter} & \textbf{Type} & \textbf{Description} \\
\hline
\texttt{file} & File & PDF, DOCX, or TXT document \\
\texttt{title} & String & Lesson title (optional, defaults to filename) \\
\texttt{subject} & String & Subject name (required) \\
\texttt{grade\_level} & String & Target grade level (required) \\
\hline
\end{tabular}
\end{table}

\textbf{Response (201 Created):}
\begin{lstlisting}
{
  "id": "f6540ed3-1020-4713-b14b-c1bc853422a8",
  "title": "Introduction to Biology",
  "subject": "Biology",
  "grade_level": "10",
  "status": "PROCESSING",
  "file_path": "uploads/f6540ed3-1020-4713-b14b-c1bc853422a8.pdf",
  "created_at": "2026-05-01T10:30:00Z"
}
\end{lstlisting}

\textbf{Status Codes:}
\begin{itemize}
    \item \textbf{201 Created}: File uploaded successfully
    \item \textbf{400 Bad Request}: Missing required parameters or invalid file type
    \item \textbf{401 Unauthorized}: Not authenticated or invalid token
    \item \textbf{403 Forbidden}: User is not a teacher or admin
    \item \textbf{500 Internal Server Error}: Server processing error
\end{itemize}

\textbf{Supported File Types:}
\begin{itemize}
    \item \texttt{.pdf} - PDF documents
    \item \texttt{.docx} - Microsoft Word documents
    \item \texttt{.txt} - Plain text files
\end{itemize}

% ============================================================================
% SECTION 5: VOICE SESSIONS ENDPOINTS
% ============================================================================
\subsection{Voice Sessions Endpoints}

Voice sessions track student interactions during learning activities.

\subsubsection{1. Start Voice Session}

\textbf{POST} \texttt{/voice/session/start}

Begin a new voice learning session for a student.

\textbf{Headers:}
\begin{lstlisting}
Authorization: Bearer <access_token>
Content-Type: application/json
\end{lstlisting}

\textbf{Request Body:}
\begin{lstlisting}
{
  "student_id": "3213b206-f779-4698-aff4-a3c3644938af",
  "note_id": "f6540ed3-1020-4713-b14b-c1bc853422a8",
  "unit_id": "308685dc-7b35-4494-b62b-f9ee2de4b884"
}
\end{lstlisting}

\textbf{Response (201 Created):}
\begin{lstlisting}
{
  "session_id": "550e8400-e29b-41d4-a716-446655440000",
  "student_id": "3213b206-f779-4698-aff4-a3c3644938af",
  "note_id": "f6540ed3-1020-4713-b14b-c1bc853422a8",
  "unit_id": "308685dc-7b35-4494-b62b-f9ee2de4b884",
  "status": "ACTIVE",
  "started_at": "2026-05-01T10:56:00Z"
}
\end{lstlisting}

\textbf{Session States:}
\begin{itemize}
    \item \textbf{ACTIVE}: Session is running
    \item \textbf{PAUSED}: Student has paused the session
    \item \textbf{COMPLETED}: Session has ended
\end{itemize}

\subsubsection{2. Log Voice Interaction}

\textbf{POST} \texttt{/voice/session/event}

Log a student action during a voice session.

\textbf{Request Body:}
\begin{lstlisting}
{
  "session_id": "550e8400-e29b-41d4-a716-446655440000",
  "interaction_type": "answer",
  "command": "The capital of France is Paris",
  "confidence": 0.95,
  "response": "Correct! Paris is indeed the capital."
}
\end{lstlisting}

\textbf{Interaction Types:}
\begin{table}[H]
\centering
\begin{tabular}{|l|l|}
\hline
\textbf{Type} & \textbf{Description} \\
\hline
\texttt{answer} & Student answers a question \\
\texttt{next} & Student requests next unit \\
\texttt{repeat} & Student requests to repeat unit \\
\texttt{free\_ask} & Student asks a free-form question \\
\texttt{pause} & Student pauses the session \\
\hline
\end{tabular}
\end{table}

\textbf{Response (201 Created):}
\begin{lstlisting}
{
  "interaction_id": "uuid-value",
  "session_id": "550e8400-e29b-41d4-a716-446655440000",
  "sequence_number": 1,
  "student_transcript": "The capital of France is Paris",
  "detected_intent": "answer",
  "ai_response_text": "Correct! Paris is indeed the capital.",
  "confidence_score": 0.95,
  "created_at": "2026-05-01T10:56:30Z"
}
\end{lstlisting}

\subsubsection{3. End Voice Session}

\textbf{POST} \texttt{/voice/session/end}

Close a voice session and calculate summary statistics.

\textbf{Request Body:}
\begin{lstlisting}
{
  "session_id": "550e8400-e29b-41d4-a716-446655440000",
  "duration_seconds": 1200,
  "questions_answered": 8,
  "total_score": 7.5
}
\end{lstlisting}

\textbf{Response (200 OK):}
\begin{lstlisting}
{
  "session_id": "550e8400-e29b-41d4-a716-446655440000",
  "status": "COMPLETED",
  "duration_seconds": 1200,
  "questions_answered": 8,
  "average_score": 0.9375,
  "ended_at": "2026-05-01T11:00:00Z"
}
\end{lstlisting}

\subsubsection{4. Get Voice Session}

\textbf{GET} \texttt{/voice/session/\{session\_id\}}

Retrieve details of a specific voice session.

\textbf{Response (200 OK):}
\begin{lstlisting}
{
  "session_id": "550e8400-e29b-41d4-a716-446655440000",
  "student_id": "3213b206-f779-4698-aff4-a3c3644938af",
  "note_id": "f6540ed3-1020-4713-b14b-c1bc853422a8",
  "unit_id": "308685dc-7b35-4494-b62b-f9ee2de4b884",
  "status": "ACTIVE",
  "started_at": "2026-05-01T10:56:00Z"
}
\end{lstlisting}

\subsubsection{5. Get Session Interactions}

\textbf{GET} \texttt{/voice/session/\{session\_id\}/interactions}

Retrieve all interactions logged during a session.

\textbf{Response (200 OK):}
\begin{lstlisting}
[
  {
    "interaction_id": "uuid-1",
    "session_id": "550e8400-e29b-41d4-a716-446655440000",
    "sequence_number": 1,
    "student_transcript": "The capital of France",
    "detected_intent": "answer",
    "ai_response_text": "Correct!",
    "confidence_score": 0.95,
    "created_at": "2026-05-01T10:56:30Z"
  },
  ...
]
\end{lstlisting}

% ============================================================================
% SECTION 6: ERROR HANDLING
% ============================================================================
\section{Error Handling}

All API errors follow a standard response format:

\begin{lstlisting}
{
  "detail": "Error message describing what went wrong"
}
\end{lstlisting}

\subsection{HTTP Status Codes}

\begin{table}[H]
\centering
\begin{tabular}{|l|l|}
\hline
\textbf{Code} & \textbf{Meaning} \\
\hline
\texttt{200} & OK - Request successful \\
\texttt{201} & Created - Resource successfully created \\
\texttt{400} & Bad Request - Invalid request parameters \\
\texttt{401} & Unauthorized - Authentication failed \\
\texttt{403} & Forbidden - Insufficient permissions \\
\texttt{404} & Not Found - Resource does not exist \\
\texttt{500} & Internal Server Error - Server-side error \\
\hline
\end{tabular}
\caption{HTTP Status Codes}
\end{table}

\subsection{Common Error Messages}

\begin{table}[H]
\centering
\begin{tabular}{|l|l|}
\hline
\textbf{Error} & \textbf{Solution} \\
\hline
\texttt{Incorrect email or password} & Verify credentials are correct \\
\hline
\texttt{Invalid or expired token} & Re-authenticate to get new token \\
\hline
\texttt{Insufficient permissions} & Check user role for endpoint access \\
\hline
\texttt{Resource not found} & Verify resource ID is correct \\
\hline
\end{tabular}
\end{table}

% ============================================================================
% SECTION 7: IMPLEMENTATION EXAMPLES
% ============================================================================
\section{Implementation Examples}

\subsection{Example 1: Complete Login Flow (Python)}

\begin{lstlisting}[language=python]
import requests
import json

API_BASE = "https://visiolearn-backend.onrender.com/api/v1"

# Step 1: Login
response = requests.post(
    f"{API_BASE}/auth/login",
    json={
        "email": "admin@visiolearn.org",
        "password": "SecurePass123!"
    }
)

if response.status_code == 200:
    tokens = response.json()
    access_token = tokens["access_token"]
    refresh_token = tokens["refresh_token"]
    print("Login successful!")
else:
    print(f"Login failed: {response.json()}")
    exit(1)

# Step 2: Use access token in subsequent requests
headers = {
    "Authorization": f"Bearer {access_token}",
    "Content-Type": "application/json"
}

# Step 3: Refresh token when expired
response = requests.post(
    f"{API_BASE}/auth/refresh",
    json={"refresh_token": refresh_token}
)

if response.status_code == 200:
    new_tokens = response.json()
    access_token = new_tokens["access_token"]
    print("Token refreshed successfully!")
\end{lstlisting}

\subsection{Example 2: Upload Lesson Note (cURL)}

\begin{lstlisting}[language=bash]
curl -X POST \
  https://visiolearn-backend.onrender.com/api/v1/notes/upload \
  -H "Authorization: Bearer YOUR_ACCESS_TOKEN" \
  -F "file=@lesson.pdf" \
  -F "title=Biology Basics" \
  -F "subject=Biology" \
  -F "grade_level=10"
\end{lstlisting}

\subsection{Example 3: Voice Session Flow (JavaScript)}

\begin{lstlisting}[language=javascript]
const API_BASE = 
  "https://visiolearn-backend.onrender.com/api/v1";
const accessToken = "YOUR_ACCESS_TOKEN";

const headers = {
  "Authorization": `Bearer ${accessToken}`,
  "Content-Type": "application/json"
};

// Start session
const startSession = async () => {
  const response = await fetch(
    `${API_BASE}/voice/session/start`,
    {
      method: "POST",
      headers,
      body: JSON.stringify({
        student_id: "3213b206-...",
        note_id: "f6540ed3-...",
        unit_id: "308685dc-..."
      })
    }
  );
  return await response.json();
};

// Log interaction
const logInteraction = async (sessionId, command) => {
  const response = await fetch(
    `${API_BASE}/voice/session/event`,
    {
      method: "POST",
      headers,
      body: JSON.stringify({
        session_id: sessionId,
        interaction_type: "answer",
        command: command,
        confidence: 0.95,
        response: "Recorded"
      })
    }
  );
  return await response.json();
};

// End session
const endSession = async (sessionId, duration) => {
  const response = await fetch(
    `${API_BASE}/voice/session/end`,
    {
      method: "POST",
      headers,
      body: JSON.stringify({
        session_id: sessionId,
        duration_seconds: duration,
        questions_answered: 8,
        total_score: 7.5
      })
    }
  );
  return await response.json();
};
\end{lstlisting}

% ============================================================================
% SECTION 8: BEST PRACTICES
% ============================================================================
\section{Best Practices}

\subsection{Authentication}

\begin{itemize}
    \item \textbf{Store tokens securely}: Never expose tokens in client-side code
    \item \textbf{Use refresh tokens}: Implement token rotation for security
    \item \textbf{Handle expiration}: Catch 401 responses and re-authenticate
    \item \textbf{Logout on close}: Revoke refresh token when user logs out
\end{itemize}

\subsection{API Requests}

\begin{itemize}
    \item \textbf{Validate input}: Check required fields before sending requests
    \item \textbf{Handle errors}: Implement proper error handling for all responses
    \item \textbf{Use pagination}: For endpoints returning large datasets
    \item \textbf{Set timeouts}: Prevent requests from hanging indefinitely
\end{itemize}

\subsection{Voice Sessions}

\begin{itemize}
    \item \textbf{Verify IDs}: Ensure student, note, and unit IDs are valid
    \item \textbf{Track state}: Maintain session state in your client
    \item \textbf{Log interactions}: Record all user actions for analytics
    \item \textbf{Handle pauses}: Manage session pauses gracefully
\end{itemize}

\subsection{File Uploads}

\begin{itemize}
    \item \textbf{Validate file type}: Check file extension before upload
    \item \textbf{Check file size}: Ensure file is not too large
    \item \textbf{Use multipart}: Use proper multipart/form-data encoding
    \item \textbf{Monitor status}: Poll endpoint to check processing status
\end{itemize}

% ============================================================================
% SECTION 9: RATE LIMITING & QUOTAS
% ============================================================================
\section{Rate Limiting \& Quotas}

\subsection{Current Limits}

\begin{itemize}
    \item \textbf{File upload size}: Maximum 50 MB
    \item \textbf{Maximum concurrent sessions}: Unlimited per user
    \item \textbf{Token expiration}: 60 minutes (access), 90 days (refresh)
\end{itemize}

\subsection{Recommendations}

\begin{itemize}
    \item Implement exponential backoff for retries
    \item Cache responses when appropriate
    \item Monitor API usage and logs regularly
    \item Contact support for quota increases
\end{itemize}

% ============================================================================
% SECTION 10: SUPPORT & RESOURCES
% ============================================================================
\section{Support \& Resources}

\subsection{Documentation}

\begin{itemize}
    \item \textbf{Interactive Docs}: \url{https://visiolearn-backend.onrender.com/docs}
    \item \textbf{ReDoc}: \url{https://visiolearn-backend.onrender.com/redoc}
    \item \textbf{OpenAPI Schema}: \url{https://visiolearn-backend.onrender.com/openapi.json}
\end{itemize}

\subsection{Deployment Information}

\begin{itemize}
    \item \textbf{Platform}: Render (render.com)
    \item \textbf{Status Page}: \url{https://status.render.com}
    \item \textbf{Response Time}: Typically < 200ms
\end{itemize}

\subsection{Contact Information}

\begin{itemize}
    \item \textbf{GitHub}: \url{https://github.com/Nyanja23/VisioLearn-Backend}
    \item \textbf{Issues}: Use GitHub Issues for bug reports
    \item \textbf{Discussions}: GitHub Discussions for questions
\end{itemize}

% ============================================================================
% APPENDIX
% ============================================================================
\appendix

\section{Quick Reference}

\subsection{Endpoint Summary}

\begin{table}[H]
\centering
\small
\begin{tabular}{|l|l|l|}
\hline
\textbf{Method} & \textbf{Endpoint} & \textbf{Purpose} \\
\hline
POST & \texttt{/auth/login} & User login \\
POST & \texttt{/auth/refresh} & Refresh access token \\
POST & \texttt{/auth/logout} & User logout \\
POST & \texttt{/notes/upload} & Upload lesson file \\
POST & \texttt{/voice/session/start} & Start voice session \\
POST & \texttt{/voice/session/event} & Log interaction \\
POST & \texttt{/voice/session/end} & End session \\
GET & \texttt{/voice/session/\{id\}} & Get session details \\
GET & \texttt{/voice/session/\{id\}/interactions} & Get interactions \\
\hline
\end{tabular}
\caption{API Endpoints Quick Reference}
\end{table}

\subsection{HTTP Headers Reference}

\begin{table}[H]
\centering
\small
\begin{tabular}{|l|l|}
\hline
\textbf{Header} & \textbf{Value} \\
\hline
\texttt{Content-Type} & \texttt{application/json} \\
\texttt{Authorization} & \texttt{Bearer <access\_token>} \\
\texttt{Accept} & \texttt{application/json} \\
\hline
\end{tabular}
\caption{HTTP Headers Reference}
\end{table}

\subsection{UUID Format}

All IDs in VisioLearn are UUIDs (Universally Unique Identifiers):

\begin{center}
\texttt{550e8400-e29b-41d4-a716-446655440000}
\end{center}

Format: \texttt{8-4-4-4-12 hexadecimal digits}

\section{Changelog}

\subsection{Version 1.0 (May 2026)}

\begin{itemize}
    \item Initial API release
    \item Authentication with JWT tokens
    \item Lesson note upload and processing
    \item Voice session tracking
    \item Student progress analytics
\end{itemize}

% ============================================================================
% END OF DOCUMENT
% ============================================================================
\end{document}
